% !TeX root = ../main.tex
% 第 7 章 Besicovitch 集与 Kakeya 集 (对应原书 pp. 95--112)

\chapter{Besicovitch 集与 Kakeya 集}\label{chap:7}

\section{引言}

Kakeya 问题有一段有趣的历史.1917 年 Besicovitch 正在研究 Riemann 积分方面的问题，遇到了如下问题：若 $f$ 是定义在平面上的 Riemann 可积函数，是否总能找到一对正交坐标轴，使 $\int f(x,y)\,\mathrm{d}x$ 对一切 $y$ 都作为 Riemann 积分存在，而且所得的 $y$ 的函数也 Riemann 可积？Besicovitch 注意到，如果能构造一个平面勒贝格测度为零、且在每一方向都含有一条线段的紧集 $F$，就会给出一个反例：因为平移 $F$（如有必要）后可设 $F$ 不含与一对固定坐标轴之一平行且有理距离的线段.令 $f$ 为由 $F$ 中至少有一个有理坐标的点构成的集合 $F_0$ 的特征函数.由于 $F$ 在每个方向都含有一条使 $F_0$ 及其补都稠密的线段，每个方向都有一条使 $f$ 不 Riemann 可积的线段.另一方面，$F$ 的不连续点集平面测度为零，故由勒贝格的著名判据，$f$ 在整个平面上 Riemann 可积.

\textcite{besicovitch1919} 成功构造了具有所需性质的集合，即如今所谓的 \newterm{besicovitch-set}{``Besicovitch 集''}{Besicovitch set}.由于当时俄国的动荡局势，他的文章流传有限，该构造后来在 \emph{Mathematische Zeitschrift}（\citeyear{besicovitch1928b}）上重新发表.

大约同一时期，\textcite{kakeya1917} 与 \textcite{fujiwarakakeya1917} 提出了求最小凸集面积的问题：在这个凸集内部可以把一条单位长线段掉头，即把它挪动到原位置但旋转 $180^{\circ}$ 而不离开该集合.他们猜测高为 $1$ 的等边三角形是这样的最小集合，并记录了 Kubota 的一个观察：若去掉凸性条件，则可以有更小的集合，即三尖内摆线.凸情形的猜测由 \textcite{pal1921} 证明，他重申了不带凸性假设的更有趣的问题.这就是后来著名的 \newterm{kakeya-problem}{Kakeya 问题}{Kakeya problem}.

在 Besicovitch 于 1924 年离开俄国后不久，人们便认识到对 Besicovitch 集作简单的修改即可给出测度任意小的 Kakeya 问题的解（\textcite{besicovitch1928b,perron1928}），该问题以出乎意料的方式得到解决.

或许其后最重要的发展再次归功于 \textcite{besicovitch1964a}，他发现了 Besicovitch 集与本书迄今所描述的几何测度论之间的基本联系.利用极对偶技术，他证明了 Besicovitch 集的存在性可以作为不规则线性可测集投影定理的一个简单的对偶结果得到.令人惊讶的是，这一对偶关系此前竟未被 Besicovitch（多年来他同时对这两个专题都感兴趣）或其他研究者注意到.

大量文章描述了这些问题的种种变式.例如，人们构造了包含所有多边形的拷贝或所有半径的圆的零测度平面集，并考虑了高维类似物的可能性.近来，这类思想已被用来回答调和分析中此前困难的问题.

值得一提的是，1958 年美国数学协会选定 Kakeya 问题作为其第一部影片的主题.任何观看这部由 Besicovitch 本人解说的迷人影片的机会都不应错过.

\section{Besicovitch 集与 Kakeya 集的构造}

本节描述 Besicovitch 集的构造，即测度为零且在每个方向都含有一条线段的集合.Besicovitch 原先的构造（\citeyear{besicovitch1919,besicovitch1928b}）已被大大简化（\textcite{perron1928,vanalphen1942,rademacher1962,schoenberg1962b,besicovitch1963c,cunningham1971,fisher1973}）.所有构造背后的基本想法是形成一个 \newterm{perron-tree}{``Perron 树''}{Perron tree}：把一个高为 $1$ 的等边三角形通过分割底边分成许多同高的小三角形，然后使这些基本三角形沿底线 $L$ 滑动不同的距离（见图~\ref{fig:7.2}）.这样的图形当然含有与 $L$ 成角至少 $60^{\circ}$ 的一切方向的单位长线段.困难在于证明用这种方式可以得到测度任意小的集合.

\begin{lemma}\label{lem:7.1}
设 $T_1$ 与 $T_2$ 是底在直线 $L$ 上的一对相邻三角形，底长为 $b$，高为 $h$.取 $\frac{1}{2}<\alpha<1$.若把 $T_2$ 沿 $L$ 滑动距离 $2(1-\alpha)b$ 使之与 $T_1$ 重叠，则所得图形 $S$ 由一个与 $T_1\cup T_2$ 位似（即相似且同样放置）的三角形 $T$（$\mathcal{L}^2(T)=\alpha^2\mathcal{L}^2(T_1\cup T_2)$）以及两个辅助三角形组成.把 $T_1\cup T_2$ 换成 $S$ 所实现的面积缩减为
\[
  \mathscr{L}^{2}(T_1\cup T_2)-\mathscr{L}^{2}(S)=\mathscr{L}^{2}(T_1\cup T_2)(1-\alpha)(3\alpha-1).
\]
（$\mathcal{L}^2$ 为平面面积即勒贝格测度.）
\end{lemma}

\begin{proof}
情形见图~\ref{fig:7.1}.利用初等几何极容易求出两图形的面积差.（注意这一计算是仿射不变的.）
\end{proof}

\begin{figure}[htbp]
  \centering
  \includegraphics[width=0.42\textwidth]{fig7-1.png}
  \caption{\cref{lem:7.1} 中的三角形滑动}
  \label{fig:7.1}
\end{figure}

\begin{theorem}\label{thm:7.2}
设 $T$ 是底在直线 $L$ 上的三角形.把 $T$ 的底分成 $2^k$ 等份，把每个分点与对顶点相连，得到 $2^k$ 个\newterm{elementary-triangle}{基本三角形}{elementary triangle} $T_1,\ldots,T_{2^k}$.只要 $k$ 取得足够大，就可以把基本三角形沿 $L$ 平移到使所得（闭）图形 $S$ 的面积任意小的位置.此外，若 $V$ 是包含 $T$ 的开集，还可以使 $S\subset V$.
\end{theorem}

\begin{proof}
这一构造需要对稍后指定的固定 $\alpha$ 值反复应用上一引理.我们先处理相邻的基本三角形对.对每个 $i$（$1\leq i\leq 2^{k-1}$），把 $T_{2i}$ 沿 $L$ 相对 $T_{2i-1}$ 移动，得到图形 $S_i^1$，它由一个三角形 $T_i^1$（与 $T_{2i-1}\cup T_{2i}$ 位似）及两个辅助三角形组成.由\cref{lem:7.1}，可以使 $\mathcal{L}^2(T_i^1)=\alpha^2\mathcal{L}^2(T_{2i-1}\cup T_{2i})$，并且（比较 $S_i^1$ 与 $T_{2i-1}\cup T_{2i}$）面积缩减了 $(1-\alpha)(3\alpha-1)\mathcal{L}^2(T_{2i-1}\cup T_{2i})$.构造的第二阶段对相邻的 $S_i^1$ 进行.对 $1\leq i\leq 2^{k-2}$，把 $S_{2i}^1$ 相对 $S_{2i-1}^1$ 平移得到 $S_i^2$.由于 $T_{2i-1}^1$ 的一边与 $T_{2i}^1$ 的对边平行且相等，由\cref{lem:7.1}，可以使 $S_i^2$ 包含一个满足 $\mathcal{L}^2(T_i^2)=\alpha^2[\mathcal{L}^2(T_{2i-1}^1)+\mathcal{L}^2(T_{2i}^1)]$ 的三角形 $T_i^2$，且 $T_{2i-1}^1$ 与 $T_{2i}^1$ 的重叠导致面积至少缩减
\[
  (1-\alpha)(3\alpha-1)[\mathcal{L}^{2}(T_{2i-1}^{1})+\mathcal{L}^{2}(T_{2i}^{1})]
  =(1-\alpha)(3\alpha-1)\alpha^{2}\mathcal{L}^{2}(T_{4i-3}\cup T_{4i-2}\cup T_{4i-1}\cup T_{4i}).
\]
如此继续；在第 $(r+1)$ 阶段，通过把 $S_{2i}^r$ 相对 $S_{2i-1}^r$ 移动（$1\leq i\leq 2^{k-r}$），使三角形 $T_{2i-1}^r$ 与 $T_{2i}^r$ 的重叠导致面积至少缩减 $(1-\alpha)(3\alpha-1)\alpha^{2r}$ 乘以移来构成 $S_{2i-1}^r$ 与 $S_{2i}^r$ 的基本三角形的总面积，从而得到 $S_i^{r+1}$.最终得到单个图形 $S_1^k$，我们就取它为 $S$.$k=3$ 时的构造阶段示于图~\ref{fig:7.2}.

比较 $\bigcup_{i}S_i^r$ 在相继的 $r$ 值下的面积，可见
\[
  \mathcal{L}^{2}(S)
  \leq\mathcal{L}^{2}(T)-(1-\alpha)(3\alpha-1)(1+\alpha^{2}+\cdots+\alpha^{2(k-1)})\mathcal{L}^{2}(T)
  =\Bigl(1-\frac{(3\alpha-1)(1-\alpha^{2k})}{(1+\alpha)}\Bigr)\mathcal{L}^{2}(T).
\]
把 $\alpha$ 取得足够接近 $1$（使 $(3\alpha-1)/(1+\alpha)$ 接近 $1$），再取足够大的 $k$，就能使 $\mathscr{L}^2(S)$ 如所需地小.

最后，若 $T$ 的底长为 $b$，则在构造中固定 $T_1$ 并相对 $T_1$ 安放其他基本三角形，任何三角形被移动的距离都不超过 $b$.于是，把原三角形分成底长至多为 $\varepsilon$ 的子三角形并对每个子三角形施行上述操作，便可以在任何基本三角形移动距离不超过 $\varepsilon$ 的前提下得到面积任意小的图形.若 $\varepsilon$ 取得足够小，该图形将含于开集 $V$ 之中.
\end{proof}

\begin{figure}[htbp]
  \centering
  \includegraphics[width=0.62\textwidth]{fig7-2.png}
  \caption{$k=3$ 时 Perron 树的构造阶段}
  \label{fig:7.2}
\end{figure}

一个有趣的问题是：当 $T$ 被分成 $m$ 个基本三角形时，确定沿底线滑动三角形所能得到的最小面积 $A_m$.Schoenberg（\citeyear{schoenberg1962a,schoenberg1962b}） 证明了 $A_m=O(1/\log m)$，但甚至不知道 $A_m$ 是否关于 $m$ 单调递减.不过，要使上述证明奏效，我们只需要 $\lim_{m\to\infty}A_m=0$.我们用\cref{thm:7.2} 来构造 Besicovitch 集.

\begin{theorem}\label{thm:7.3}
存在勒贝格测度为零且在每个方向都含有一条单位长线段的平面集合.
\end{theorem}

\begin{proof}
我们构造一个测度为零、在一个 $60^{\circ}$ 扇形内的一切方向都含有单位长线段的集合 $F$.把 $F$ 与其旋转 $60^{\circ}$ 与 $120^{\circ}$ 的全等拷贝取并，便得到具有所需性质的集合.

构造由反复应用\cref{thm:7.2} 组成.设 $S_1$ 是基于直线 $L$ 的高为 $1$ 的等边三角形.取包含 $S_1$ 的开集 $V_1$ 使 $\mathcal{L}^2(\overline{V}_1)\leq 2\mathcal{L}^2(S_1)$（横线表示闭包）.由\cref{thm:7.2}，可以把 $S_1$ 分成若干个高为 $1$、底在 $L$ 上的基本三角形，并把它们沿 $L$ 平移，形成含于 $V_1$ 内且 $\mathcal{L}^2(S_2)\leq 2^{-2}$ 的闭图形 $S_2$.由于 $S_2$ 是三角形的有限并，可找到开集 $V_2$ 使 $S_2\subset V_2\subset V_1$ 且 $\mathcal{L}^2(\overline{V}_2)\leq 2\mathcal{L}^2(S_2)$.类似地，可以把 $S_2$ 的每个基本三角形进一步分成基本三角形，并沿 $L$ 移动它们使所得图形 $S_3$ 含于 $V_2$ 且 $\mathcal{L}^2(S_3)\leq 2^{-3}$.把 $S_3$ 包进满足 $\mathcal{L}^2(\overline{V}_3)\leq 2\mathcal{L}^2(S_3)$ 的开集 $V_3$，如此继续.我们得到图形序列 $\{S_i\}$——每个都是基于 $L$、高为 $1$ 的基本三角形的有限并——以及开集序列 $\{V_i\}$，使得对每个 $i$，
\[
  S_i\subset V_i\subset V_{i-1}
\]
且
\[
  \mathcal{L}^{2}(\bar{V}_i)\leq 2\mathcal{L}^{2}(S_i)\leq 2^{-i+1}.
\]
令 $F$ 为闭集 $F=\bigcap_{i=1}^{\infty}\overline{V}_i$.我们断言 $F$ 具有所需的性质.当然 $\mathscr{L}^2(F)=0$.由构造，每个 $S_i$（从而每个 $\overline{V}_i$）在与 $L$ 成角 $60^{\circ}$ 或更大的任何方向上含有单位长线段.必须证明这对 $F$ 也成立.设 $\theta$ 是这样一个方向，并对每个 $i$ 取方向 $\theta$ 的含于 $\overline{V}_i$ 的单位长线段 $M_i$.由标准的紧性论证，必要时取子列，可设 $\{M_i\}$ 按显然的意义收敛于方向 $\theta$ 的单位长线段 $M$.由于 $\{\overline{V}_i\}$ 递减，当 $i\geq j$ 时 $M_i\subset\overline{V}_j$，故因 $\overline{V}_j$ 闭，对每个 $j$ 有 $M\subset\overline{V}_j$.于是 $M\subset\bigcap_{i=1}^{\infty}\overline{V}_i=F$，如所需.
\end{proof}

人们还找到了构造 Besicovitch 集的其他方法：\textcite{kahane1969} 注意到把 $x$ 轴上一个 Cantor 集的点与一条平行的 Cantor 型集的点连接起来便能得到这样的集合，另见 \textcite{kinney1970} 与 \textcite{alexander1975}.另一种使用对偶性的方法将在 7.3~节描述.

同样的构造可以用来形成 \newterm{nikodym-set}{Nikodym 集}{Nikodym set}：即单位正方形的子集 $F$，$\mathscr{L}^2(F)=1$，使得对每个 $x\in F$ 都有一条过 $x$ 且与 $F$ 仅交于点 $x$ 的直线（我们称 $F$ 的每一点都是\newterm{linearly-accessible-set}{``线性可达的''}{linearly accessible set}）.\textcite{nikodym1927} 为回答 Banach 的一个问题给出的复杂构造被 \textcite{davies1952a} 简化，后者证明甚至可以找到这样的集合：过其每一点都有不可数多条这样的直线.（这些构造在 de Guzmán（\citeyear{deguzman1975,deguzman1979,deguzman1981}）、\textcite{casas1978} 与 \textcite{casasdeguzman1981} 中也有描述.）\textcite{davies1952a} 用类似想法得到下述令人惊讶的结果，此处引用而不给出证明.

\begin{theorem}\label{thm:7.4}
给定任何具有有限 $\mathscr{L}^2$-测度的平面集 $G$，存在由直线的并表示的集合 $F$ 使 $G\subset F$ 且 $\mathscr{L}^2(G)=\mathscr{L}^2(F)$.
\end{theorem}

\textcite{larman1971a} 也在其\newterm{impossible-set}{``不可能集''}{impossible set}的构造中使用了这类技术.他给出了一个由 $\mathbb{R}^n$（$n\geq 3$）中闭线段的不交并构成的紧集 $F$，使 $\mathscr{L}^n(F)>0$ 而 $\mathscr{L}^n(G)=0$，其中 $G$ 是相应开线段之并.于是 $F$ 的测度集中于线段的端点.（这样的构造在 $\mathbb{R}^2$ 中是不可能的.）

接下来我们用\cref{thm:7.2} 解决 Kakeya 问题.需要下述简单引理.

\begin{lemma}\label{lem:7.5}
设 $L_1$ 与 $L_2$ 是平面中的平行直线.则给定 $\varepsilon>0$，存在包含 $L_1$ 与 $L_2$ 且 $\mathscr{L}^2(E)<\varepsilon$ 的集合 $E$，使一条单位长线段可以连续地从 $L_1$ 移动到 $L_2$ 而不离开 $E$.
\end{lemma}

\begin{proof}
设 $x_1$ 与 $x_2$ 分别是 $L_1$ 与 $L_2$ 上的点.令 $E$ 为由 $L_1$、$L_2$、连接 $x_1$ 与 $x_2$ 的线段 $M$ 以及以 $x_i$ 为心、介于 $L_i$ 与 $M$ 之间的单位扇形（$i=1,2$）构成的集合，见图~\ref{fig:7.3}.把 $x_1$ 与 $x_2$ 取得相距足够远，$E$ 的总面积便可任意小.进一步，一条单位长线段可以通过在第一个扇形中旋转、沿 $M$ 平移、再在第二个扇形中旋转而从 $L_1$ 移动到 $L_2$.
\end{proof}

\begin{figure}[htbp]
  \centering
  \includegraphics[width=0.5\textwidth]{fig7-3.png}
  \caption{\cref{lem:7.5} 中的集合 $E$}
  \label{fig:7.3}
\end{figure}

\begin{theorem}\label{thm:7.6}
给定 $\varepsilon>0$，存在满足 $\mathscr{L}^2(E)<\varepsilon$ 的\newterm{kakeya-set}{Kakeya 集}{Kakeya set} $E$，在其中一条单位长线段可以被连续移动到原位置但旋转过 $180^{\circ}$.
\end{theorem}

\begin{proof}
我们构造满足 $\mathscr{L}^2(E_0)<\frac{1}{3}\varepsilon$ 的集合 $E_0$，在其内部可以把一条单位长线段挪动到与其原方向成 $60^{\circ}$ 角的位置.取这一构造的三份拷贝便可得到所需的集合 $E$.（若需要在三个组成部分之间作平移，可能要用到\cref{lem:7.5}.）

若 $T$ 是基于直线 $L$ 的高为 $1$ 的等边三角形，则可把 $T$ 分成 $m=2^k$ 个基本三角形并沿 $L$ 把它们滑动到位置 $T_1',\ldots,T_m'$，使 $\mathscr{L}^2(\bigcup_1^m T_i')<\varepsilon/6$（见\cref{thm:7.2}）.令 $E_1=\bigcup_1^m T_i'$.注意对每个 $i$，$T_i'$ 的一边与 $T_{i+1}'$ 的对边平行.于是由\cref{lem:7.5}，对每个 $i$ 可以向 $E_1$ 添加一个测度至多为 $\varepsilon/6m$ 的集合，使一条单位长线段能从 $T_i'$ 移动到 $T_{i+1}'$.这就得到测度至多为 $\frac{1}{6}\varepsilon+\frac{1}{6}(m-1)\varepsilon/m<\frac{1}{2}\varepsilon$ 的集合 $E_0$，在其中一条单位长线段可旋转过 $60^{\circ}$，如所需.
\end{proof}

\section{对偶方法}

本节把 Besicovitch 集与\cref{chap:3,chap:6} 所述的线性可测集理论联系起来.基本想法是用点来参数化直线，使得点集 $E$ 在某方向上的投影在几何上相似于由 $E$ 参数化的直线与某条固定直线的交.于是可以通过把 $E$ 沿各个方向投影来考察由这些直线之并构成的集合（``直线集''）.

\textcite{besicovitch1964a} 以这种方式使用了\newterm{polar-reciprocity}{极对偶}{polar reciprocity}技术.设 $L(x)$ 为与原点距离 $1/|x|$ 且垂直于径向量 $x$ 的直线；即 $L(x)$ 是 $x$ 关于单位圆 $C$ 的极线.无穷远点按自然方式表示过原点的直线.极对偶的基本性质是：$x\in L(y)$ 当且仅当 $y\in L(x)$.容易看出，若 $\operatorname{proj}_{\theta}$ 表示到过原点、方向 $\theta$ 的直线 $L_\theta$ 上的投影，则对任何集合 $E$，直线集 $\{L(x):x\in E\}$ 与 $L_\theta$ 的交恰是集合 $\operatorname{proj}_{\theta}E$ 关于 $C$ 的几何反演.因此\cref{chap:6} 的投影定理可用于研究直线族的交的性质.

出于技术上的原因，我们引入\newterm{duality}{对偶性}{duality}的另一种形式，尽管基本思想相同.暂时我们在平面内工作，以 $(x,y)$ 等表示 Descartes 坐标.

若 $(a,b)\in\mathbb{R}^2$，令 $L(a,b)$ 表示直线 $y=a+bx$ 上的点之集.若 $E$ 是 $\mathbb{R}^2$ 的子集，令 $L(E)$ 为直线集 $\bigcup_{(a,b)\in E}L(a,b)$.设 $c$ 为常数，$L_c$ 为直线 $x=c$，则
\[
  L(a,b)\cap L_c=(c,\,a+bc)=(c,\,(a,b)\cdot(1,c)),
\]
其中 ``$\cdot$'' 表示 $\mathbb{R}^2$ 中通常的数量积.于是若 $E\subset\mathbb{R}^2$，
\[
  L(E)\cap L_c=\{(c,\,(a,b)\cdot(1,c)):(a,b)\in E\},
\]
故集合 $L(E)\cap L_c$ 与 $\operatorname{proj}_{\theta}E$ 几何相似，相似比为 $(1+c^2)^{\frac{1}{2}}$，其中 $c=\tan\theta$.（照例，$\operatorname{proj}_{\theta}$ 表示到过原点、与 $x$ 轴成角 $\theta$ 的直线 $L_\theta$ 上的投影.）特别地，
\begin{equation}\label{eq:7.1}
  \mathscr{L}^{1}\bigl(L(E)\cap L_c\bigr)=0
  \quad\Longleftrightarrow\quad
  \mathscr{L}^{1}(\operatorname{proj}_{\theta}E)=0.
\end{equation}
进一步，考虑到 $y$ 轴上的投影可见：若 $(a,b)\in\mathbb{R}^2$，则 $\operatorname{proj}_{\pi/2}(a,b)=b$ 恰是直线 $L(a,b)$ 的斜率.于是若 $E$ 是任意点集且 $b\in\operatorname{proj}_{\pi/2}E$，则 $L(E)$ 含有一条斜率为 $b$ 的直线.

下面的定理给出线性可测集与直线集之间的基本关系.

\begin{theorem}\label{thm:7.7}
设 $E$ 是 $\mathbb{R}^2$ 中的 $1$-集.则 $L(E)$ 是 $\mathbb{R}^2$ 的 $\mathscr{L}^2$-可测子集.若 $E$ 不规则，则 $\mathscr{L}^2(L(E))=0$；若 $E$ 正则，则 $\mathscr{L}^2(L(E))>0$.
\end{theorem}

\begin{proof}
唯一稍显麻烦的是验证 $L(E)$ 的可测性，不过照例这并不构成严重的困难（见 \textcite{croft1965} 在其对 Besicovitch 文章的评论中的说明）.

容易看出，若 $E$ 是开集则 $L(E)$ 是开集，若 $E$ 是闭集则 $L(E)$ 是闭集.于是若 $E$ 是 $G_{\delta}$-集（开集的可数交）或 $F_{\sigma}$-集（闭集的可数并），则 $L(E)$ 也是同一类型的集合.

设 $\mathcal{H}^1(E)=0$.由 $\mathcal{H}^1$ 的正则性（\cref{thm:1.6}），可找到 $G_{\delta}$-集 $E_0$ 使 $E\subset E_0$ 且 $\mathcal{H}^1(E_0)=0$.由\cref{lem:6.1}，对一切 $\theta$ 有 $\mathscr{L}^1(\operatorname{proj}_{\theta}E_0)=0$，故由 \eqref{eq:7.1}，对一切 $c$ 有 $\mathscr{L}^1(L(E_0)\cap L_c)=0$.但 $L(E_0)$ 是 $G_{\delta}$-集，从而是平面可测的.Fubini 定理适用，故 $\mathscr{L}^2(L(E_0))=0$，由此得 $L(E)\subset L(E_0)$ 是 $\mathscr{L}^2$-可测的且 $\mathscr{L}^2(L(E))=0$.

现在设 $E$ 是任意 $1$-集.由正则性，$E=E_0\cup F$，其中 $E_0$ 是 $F_{\sigma}$-集且 $\mathcal{H}^1(F)=0$.于是 $L(E)=L(E_0)\cup L(F)$ 作为两个可测集之并是 $\mathscr{L}^2$-可测的.

若 $E$ 是不规则 $1$-集，则由\cref{thm:6.13}，对几乎所有的 $\theta$ 有 $\mathscr{L}^1(\operatorname{proj}_{\theta}E)=0$.由对偶原理 \eqref{eq:7.1}，对几乎所有的 $c$ 有 $\mathscr{L}^1\bigl(L(E)\cap L_c\bigr)=0$.由于 $L(E)$ 是 $\mathscr{L}^2$-可测的，Fubini 定理蕴含 $\mathscr{L}^2(L(E))=0$.

若 $E$ 是正则 $1$-集，则由\cref{thm:6.10}，对几乎所有的 $\theta$ 有 $\mathscr{L}^1(\operatorname{proj}_{\theta}E)>0$，故由 \eqref{eq:7.1}，对几乎所有的 $c$ 有 $\mathscr{L}^1\bigl(L(E)\cap L_c\bigr)>0$.此时由 $L(E)$ 的可测性与 Fubini 定理得 $\mathscr{L}^2(L(E))>0$.
\end{proof}

\textcite{besicovitch1964a} 进一步证明：若 $E$ 是正则 $1$-集，则只要把直线集覆盖的重数考虑在内，$L(E)$ 就具有无穷平面测度.然而 \textcite{davies1965} 构造了一个满足 $\mathscr{L}^2(L(E))<\infty$ 的正则 $1$-集 $E$，表明重数假设是关键的.（虽然这些文章用极对偶表示直线，不同的参数化在一个局部实解析双射下等价，因此集合的正则性与所用的参数化无关.）

现在我们用对偶性给出 Besicovitch 集的一个简单构造.这加强了\cref{thm:7.3}，因为我们得到的是在每个方向都含有一条完整直线（而不仅仅是一条线段）的零测度集合.

\begin{theorem}\label{thm:7.8}
存在平面中勒贝格测度为零且在每个方向都含有一条直线的子集.
\end{theorem}

\begin{proof}
设 $E$ 是 $\mathbb{R}^2$ 中的不规则 $1$-集，且 $E$ 到 $y$ 轴上的投影 $\operatorname{proj}_{\pi/2}E$ 包含线段 $-1\leq y\leq 1$.例如，$E$ 可以取为\cref{thm:3.32} 的集合连同其关于 $x$ 轴的反射.考虑直线集 $L(E)$；由\cref{thm:7.7}，$\mathscr{L}^2(L(E))=0$.另一方面，\cref{thm:7.7} 之前的注记表明 $L(E)$ 含有斜率介于 $-1$ 与 $1$ 之间的一切直线.把 $L(E)$ 与其旋转 $\frac{1}{2}\pi$ 的全等拷贝取并，便得到具有所需性质的 $\mathbb{R}^2$ 的子集.
\end{proof}

我们对偶性的下一个应用归功于 \textcite{davies1971}，它表明 Besicovitch 集必然相当大.

\begin{theorem}\label{thm:7.9}
设 $F$ 是平面中在每个方向都含有一条直线的子集.则 $F$ 的 Hausdorff 维数为 $2$.
\end{theorem}

\begin{proof}
任何集合都含于一个同维数的 $G_{\delta}$-集，故不失一般性可设 $F$ 是 $G_{\delta}$-集.于是集合
\[
  E=\{(a,b):L(a,b)\subset F\}
  =\bigcap_{r=1}^{\infty}\{(a,b):L(a,b)\cap B_r(0)\subset F\cap B_r(0)\}
\]
也是 $G_{\delta}$-集，从而是 Borel 可测的.由于 $F$ 在每个方向都含直线，$\operatorname{proj}_{\pi/2}E$ 是整条 $y$ 轴，故由\cref{lem:6.1} 有 $\mathcal{H}^1(E)=\infty$.由投影定理（\cref{thm:6.8}(a)），对几乎所有的 $\theta$ 有 $\dim(\operatorname{proj}_{\theta}E)=1$，故由对偶原理，对几乎所有的 $c$ 有 $\dim\bigl(L(E)\cap L_c\bigr)=1$.由\cref{thm:5.8}，$\dim L(E)=2$，于是由 $L(E)\subset F$ 得 $\dim F=2$.
\end{proof}

把上述证明稍作修改可以证明：若 $F$ 在每个方向含有一条线段，则 $\dim F=2$（见 \textcite{davies1971}）.事实上还有更强的结果：若 $\dim F<2$，则 $F$ 在几乎所有的方向上与一切直线交成 $1$ 维测度为零的集合.这可以用相同的基本思想并采用 6.3~节末尾提到的投影定理的 Marstrand 变式来证明.（这一技术在\citeyear{falconer1980b} 中被用于高维.）

\textcite{besicovitchrado1968}（在 Besicovitch 的最后一篇文章中）与 \textcite{kinney1968} 考虑了另一个曲线填装问题，给出了包含每个半径的圆的零测度平面集的直接构造.而正是 \textcite{davies1972} 认识到这一问题适合对偶方法.

\begin{theorem}\label{thm:7.10}
存在平面中勒贝格测度为零且包含每个半径的圆的子集.
\end{theorem}

\begin{proof}
设 $E$ 是一个有界不规则 $1$-集（例如\cref{thm:3.32} 所描述者），其到 $x$ 轴上的投影 $\operatorname{proj}_{\theta}E$ 包含区间 $0\leq x\leq 1$.如下定义映射 $\psi:\mathbb{R}^2\to\mathbb{R}^2$：
\[
  \psi(x,y)=\bigl(x(1+y^{2})^{\frac{1}{2}},\,y\bigr).
\]
对每个满足 $0\leq x\leq 1$ 的 $x$，存在 $y$ 使 $(x,y)\in E$.于是对任何这样的 $x$ 值，可以在 $\psi(E)$ 中找到形如 $\bigl(x(1+y^2)^{\frac{1}{2}},y\bigr)$ 的点，即找到 $(a,b)\in\psi(E)$ 使 $a^2-x^2b^2=x^2$ 或 $x=|a|(1+b^2)^{-\frac{1}{2}}$.而这恰是直线 $L(a,b)$ 到原点的垂直距离，故直线集 $L(\psi(E))$ 含有与原点距离介于 $0$ 与 $1$ 之间的一切直线.进一步，由于 $\psi$ 是实解析映射，$\psi(E)$ 是不规则 $1$-集（$\psi$ 把可求长曲线映为可求长曲线），故由\cref{thm:7.7}，$\mathscr{L}^2(L(\psi(E)))=0$.现在令 $F$ 为把 $L(\psi(E))$ 关于原点反演所得的集合（换言之，$F$ 是 $L(\psi(E))$ 在由极坐标 $(r,\theta)\to(1/r,\theta)$ 给出的变换下的像）.反演把零测度集变为零测度集，并把距原点 $x$ 的直线变为过原点、半径为 $1/(2x)$ 的圆.于是 $F$ 测度为零，且含有半径大于 $\frac{1}{2}$ 的一切圆.取与 $F$ 相似、相似比趋于零的集合的可数并，便得到具有所需性质的集合.
\end{proof}

鉴于这一构造，自然会问：测度为零的平面集是否可能包含圆心位于每一点的圆.这个有趣而困难的问题至今未解.

Davies 的定理（\cref{thm:7.4}）有一个令人惊讶的对偶.实质上，我们可以找到一个集合，使它在每个方向上具有任何想要的投影或``影子''.这推广了 \textcite{talagrand1980} 的结果.

\begin{theorem}\label{thm:7.11}
对每个 $\theta\in[0,\pi)$ 设 $G_\theta$ 是 $\mathbb{R}$ 的子集，并设集合 $\{(\theta,y):y\in G_\theta\}$ 是平面勒贝格可测的.则存在集合 $E\subset\mathbb{R}^2$，使得对几乎所有的 $\theta$ 有 $\{u\boldsymbol{\theta}:u\in G_\theta\}\subset\operatorname{proj}_{\theta}E$ 且 $\mathscr{L}^1(G_\theta)=\mathscr{L}^1(\operatorname{proj}_{\theta}E)$，其中 $\boldsymbol{\theta}$ 是方向 $\theta$ 的单位向量.
\end{theorem}

\begin{proof}
令
\begin{equation}\label{eq:7.2}
  G=\bigl\{(c,y):(1+c^{2})^{-\frac{1}{2}}y\in G_{\theta},\ \text{其中}\ c=\tan\theta\bigr\}.
\end{equation}
由\cref{thm:7.4}，可找到直线集（记作 $L(E)$）包含 $G$ 且 $\mathscr{L}^2(G)=\mathscr{L}^2(L(E))$.由于 $G$ 平面可测，由 Fubini 定理，对几乎所有的 $c$ 有 $\mathscr{L}^1(G\cap L_c)=\mathscr{L}^1\bigl(L(E)\cap L_c\bigr)$，且 $G\cap L_c\subset L(E)\cap L_c$.由对偶原理，集合 $L(E)\cap L_c$ 与 $\operatorname{proj}_{\theta}E$ 相似（$c=\tan\theta$），相似比为 $(1+c^2)^{\frac{1}{2}}$.于是由 \eqref{eq:7.2}，对几乎所有的 $\theta$，$\operatorname{proj}_{\theta}E$ 与 $G_\theta$ 至多相差一个零测集.
\end{proof}

最后，我们用对偶性考察 Besicovitch 集的高维类似物的可能性.

\begin{theorem}\label{thm:7.12}
$\mathbb{R}^3$ 中 $\mathscr{L}^3$-测度为零的子集不可能包含每个平面的一个平移.
\end{theorem}

\begin{proof}
对 $\mathbb{R}^3$ 中的平面作适当的参数化后，证明与\cref{thm:7.9} 的证明非常相似.对 $(a,b,c)\in\mathbb{R}^3$，令 $\Pi(a,b,c)$ 表示平面 $z=a+bx+cy$ 上的点之集；若 $E\subset\mathbb{R}^3$，令 $\Pi(E)=\bigcup_{(a,b,c)\in E}\Pi(a,b,c)$.若 $d,e$ 为常数，$L_{d,e}$ 为直线 $x=d,\ y=e$，则
\[
  \Pi(a,b,c)\cap L_{d,e}=(d,e,\,a+bd+ce)=(d,e,\,(a,b,c)\cdot(1,d,e)),
\]
其中 ``$\cdot$'' 是 $\mathbb{R}^3$ 中的数量积.于是若 $E\subset\mathbb{R}^3$，
\[
  \Pi(E)\cap L_{d,e}=\{(d,e,\,(a,b,c)\cdot(1,d,e)):(a,b,c)\in E\},
\]
故 $\Pi(E)\cap L_{d,e}$ 与 $\operatorname{proj}_{\theta}E$ 几何相似，后者是 $E$ 到方向为 $(1,d,e)$ 的直线上的投影.于是
\begin{equation}\label{eq:7.3}
  \mathscr{L}^{1}\bigl(\Pi(E)\cap L_{d,e}\bigr)=0
  \quad\Longleftrightarrow\quad
  \mathscr{L}^{1}(\operatorname{proj}_{\theta}E)=0.
\end{equation}
现在设集合 $F\subset\mathbb{R}^3$ 包含每个平面的一个平移.不失一般性可设 $F$ 是 $G_{\delta}$-集，此时容易看出集合 $E=\{(a,b,c):\Pi(a,b,c)\subset F\}$ 是 Borel 集.若 $F$ 包含与所有单位向量垂直的平面，且 $\Pi$ 为平面 $x=0$，则 $\operatorname{proj}_{\Pi}E$ 是整个 $\Pi$，故由\cref{lem:6.1}，$\mathcal{H}^2(E)=\infty$.由\cref{thm:6.9}(b)，对几乎所有的 $\theta$ 有 $\mathscr{L}^1(\operatorname{proj}_{\theta}E)>0$，从而由 \eqref{eq:7.3}，对几乎所有的 $(d,e)\in\mathbb{R}^2$ 有 $\mathscr{L}^1\bigl(\Pi(E)\cap L_{d,e}\bigr)>0$.与\cref{thm:7.9} 的证明一样，$E$ 的可测性蕴含 $\Pi(E)$ 是 $\mathscr{L}^3$-可测的，故由 Fubini 定理，$\mathscr{L}^3(\Pi(E))>0$.由于 $\Pi(E)\subset F$，集合 $F$ 必有正测度.
\end{proof}

把上述证明稍作修改并不难证明：$\mathbb{R}^3$ 中任何零测度子集与几乎所有方向（按球面测度）垂直的一切平面都交成平面测度为零的集合.

Besicovitch 集的高维类似物是否存在是颇有意义的问题.称 $\mathbb{R}^n$ 的子集为 $(n,k)$-Besicovitch 集，如果它 $n$ 维勒贝格测度为零，却包含 $\mathbb{R}^n$ 的每个 $k$ 维子空间的一个平移.我们在\cref{thm:7.8} 中证明了 $(2,1)$-Besicovitch 集存在，而且对任意 $n$，这样的集合与 $\mathbb{R}^{n-2}$ 的 Descartes 积是一个 $(n,1)$-Besicovitch 集.\textcite{marstrand1979b} 用直接的覆盖方法证明了 $(3,2)$-集不可能存在；同时 \textcite{falconer1980a} 用 Fourier 变换证明：当 $k>\frac{1}{2}n$ 时 $(n,k)$-Besicovitch 集不存在.\textcite{oberlinstein1982} 也联系泛函分析提到了 $(n,n-1)$ 情形.\cref{thm:7.12} 足以推广到 $(n,n-1)$-Besicovitch 集不存在的结论，但在一般的 $(n,k)$ 情形使用对偶比较麻烦.\textcite{falconer1980b} 通过使用一种对偶形式并结合投影定理（投影经由一个公共的中间空间分解），证明了当 $2\leq k\leq n-1$ 时 $(n,k)$-Besicovitch 集不存在.这里我们陈述这一结果的最一般形式：

\begin{theorem}\label{thm:7.13}
设 $2\leq k\leq n-1$，$F$ 是 $\mathbb{R}^n$ 中 $n$ 维勒贝格测度为零的子集.则对几乎所有的 $k$ 维子空间 $\Pi\subset\mathbb{R}^n$（按 Grassmann 流形 $G_{n,k}$ 上通常的不变测度），$\Pi$ 的每个平移与 $F$ 交成 $k$ 维测度为零的集合.
\end{theorem}

\section{推广}

本节综述由 Besicovitch 与 Kakeya 构造引发的众多推广与变式.（其中一些材料在 Croft 未刊的油印讲义以及 \textcite{davies1971,cunningham1974,taylor1975} 中也有讨论.）

一类一般形式的问题是：给定一族几何图形，是否可能找到一个零测度平面集，包含该族中每个图形的一个平移，或哪怕只是一个全等拷贝？Besicovitch 集对全体直线的平移给出了肯定答案.

\textcite{kinney1968} 观察到，由于 Cantor 集 $E$ 勒贝格测度为零且包含 $(0,1)$ 中的一切距离，``Cantor 格子布''$(E\times[0,1])\cap([0,1]\times E)$ 平面测度为零，且包含边长小于 $1$ 的每个矩形的全等拷贝.取相似集合的可数并便得到包含所有矩形拷贝的零测度集.Ward（\citeyear{ward1970a,ward1970b}） 构造了包含所有多边形的全等拷贝的零测度集；事实上他造出了一个 Hausdorff 维数为 $2-1/n$ 的集合，包含每个 $n$ 边多边形的一个全等拷贝.\textcite{davies1971} 随后修改了 Besicovitch 的原始构造，给出包含每条多边形折线的平移的零测度集——由\cref{thm:7.9} 它必有维数 $2$.\textcite{marstrand1972} 在一篇适用于任何拓扑测度空间的非常一般的文章中证明了这类最强的结果.他证明了，除其他事项外，存在包含每条直线的可数并的平移的零测度集，还存在 Hausdorff 维数小至 $1$、包含每个这样的并的全等拷贝的集合.

另一个自然会考虑的问题是\newterm{curve-packing}{曲线填装}{curve packing}：把曲线弧的拷贝填装进零测度集.我们已经见过（\cref{thm:7.10} 及相关注记）存在包含所有半径的圆的零测度集，也提到了所有圆心的问题尚未解决.其他显然的问题是：零测度集能否包含所有椭圆、所有二次曲线、或每个度数 $k>1$ 的平面曲线的拷贝.虽然人们普遍相信这样的集合不存在，这些问题相当困难，部分原因是其中涉及非线性，迄今一无所证.\textcite{falconer1982} 给出了关于此类集合的维数以及在 $\mathbb{R}^n$（$n\geq 3$）中填装高维曲面的一些相对较弱的结果.其他这样的填装结果（通常在三维或更高维）可以从调和分析的极大定理推出，见 7.5~节.

毫无疑问，这一领域最重要的发展归功于 \textcite{marstrand1979a}，他解决了\newterm{worm-problem}{``虫''问题}{worm problem}（``一条单位长的虫最小的舒适住所是什么？''），证明了对 $\mathbb{R}^n$（$n\geq 2$）中任何 $n$ 维测度为零的子集 $F$，存在某条长度为 $1$ 的光滑（$C^\infty$）曲线，其全等拷贝都不含于 $F$.事实上，存在这样的曲线，它不能被任何在包含该曲线的区域内定义的可逆实解析映射变换成 $F$ 的子集.复杂的证明使用了全有界度量空间的熵的思想，即对小的 $\varepsilon$ 估计覆盖该空间所需的直径为 $\varepsilon$ 的集合的个数.

还应简要提及相反类型的填装问题：是否存在不包含某些指定集合的全等或相似拷贝的正测度集合？这类结果中最著名的归功于 \textcite{steinhaus1920}：直线的任何正测度子集包含 $(0,c)$ 中的一切距离，其中 $c$ 为某个正数.（见\cref{exe:1.7}；各种推广见 \textcites{besicovitchmiller1948,besicovitch1948,eggleston1949,besicovitchtaylor1952}.）测度为零的 Cantor 集也具有这一性质.Steinhaus 结果由勒贝格密度定理（\cref{thm:1.13}）推出，后者还蕴含直线的任何正长度子集包含任何有限点集的相似拷贝.Erdős 提出了诱人的问题：这对可数（收敛）点列是否仍成立；\textcite{falconer1984} 给出了部分解答.\textcite{daviesmarstrandtaylor1960} 与 \textcite{darstgoffman1970} 考察了这一主题的其他若干变式.

求包含某些类型图形的所有拷贝的最小凸集的问题实际上属于凸性而非几何测度论.我们仅指出：在激发本章所描述的许多研究的 \textcite{pal1921} 的文章中，证明了等边三角形是包含每个方向的单位长线段的最小平面凸集；高维的类似问题仍未解决.\textcite{eggleston1957,besicovitch1965b,wetzel1973} 等人考虑了相关的``凸性''问题.

我们回到在一个尽可能小的集合内把单位长线段转动 $180^{\circ}$ 的 Kakeya 问题.最终由 \textcite{cunningham1971} 解决，他构造了含于单位圆盘内的测度任意小的单连通 Kakeya 集.同一篇文章考虑了极小星形 Kakeya 集，并证明其面积的下界介于 $\pi/108$ 与 $(5-2\sqrt{2})\pi/24$ 之间，取代了 \textcite{walker1952,blank1963,cunninghamschoenberg1965} 更早的估计.（称集合是星形的，如果它有一点可用含于该集合内的线段连接到集合的所有其他点.）极小凸 Kakeya 集仍是高为 $1$ 的等边三角形（\citeyear{pal1921}），三维类似物未知.

\textcite{cunningham1974} 考虑了值得一提的三个问题.首先他求出球面上可以在其中把一段大圆弧掉头的最小子集的面积.极小面积（对小于半圆的弧，此前已由 \textcite{wilker1971} 证明为零）取决于弧所张的角.文中讨论的第二个问题涉及求出小的（平面）集合，使一条定角而半径可变的圆弧可以在其中收缩成一点.Cunningham 的最后一个变式把``鸟''定义为一条中央线段（``身体'')，其两端各非刚性地连接一段（``翅膀''）.文中证明可以把鸟连续地移过任何有界集 $E$，使身体经过 $E$ 的全部，而翅膀保持在一个测度任意小的集合内.这一问题在调和分析中有应用，并与 \textcite{davies1952a} 关于线性可达集的工作密切相关.

\section{与调和分析的关系}

若不提及与调和分析的关系，本章便不完整.一方面，有时可以从某些积分的更一般的范数界得到填装结果.另一方面，Kakeya 构造被用来为调和分析中的一些重大猜想提供反例.我们宽泛地使用\newterm{harmonic-analysis}{``调和分析''}{harmonic analysis}一词，用以描述包括 Fourier 分析以及微分理论、极大算子与乘子理论在内的广阔数学领域.我们只能给出少量例子来说明这种关系.

在第一个例子中，我们从泛函分析不等式推出\cref{thm:7.12}（高维 Besicovitch 集的不存在性）.设 $f(x)$ 是定义在 $\mathbb{R}^3$ 上的可测函数.照例记 $\|f\|_p=\bigl[\int_{\mathbb{R}^3}|f|^p\,\mathrm{d}x\bigr]^{1/p}$，$1\leq p<\infty$.令 $F(t,\theta)$ 表示 $f$ 在垂直于单位向量 $\theta$、距原点距离为 $t$ 的平面上的积分；由 Fubini 定理的简单应用，对一切 $\theta$，$F(t,\theta)$ 对几乎所有的 $t$ 存在.把 $f$ 的 Fourier 变换关于 $t$ 沿所有单位向量 $\theta$ 积分，并利用可积函数的变换有界这一事实，不难证明
\begin{equation}\label{eq:7.4}
  \int_{\theta}\operatorname*{ess\,sup}_{t}|F(t,\theta)|\,\mathrm{d}\theta\leq c\bigl(\|f\|_{1}+\|f\|_{2}\bigr)
\end{equation}
对某个与 $f$ 无关的常数 $c$ 成立.（见 \textcite{falconer1980a} 或 \textcite{oberlinstein1982}.）

现在设 $E$ 是零测度集，并记 $F_0(\theta)$ 为 $E$ 与垂直于 $\theta$ 的平面的交的平面外测度的上确界.常规方法可证 $F_0$ 关于球面测度可测.由勒贝格测度的正则性，可找到包含 $E$ 且测度小于 $\varepsilon$ 的开集 $V$；于是若 $f$ 是 $V$ 的特征函数，则 $\|f\|_1=\|f\|_2<\varepsilon$.由于 $V$ 开，
\[
  F_0(\theta)\leq\operatorname*{ess\,sup}_{t}F(t,\theta),
  \quad\text{故由 \eqref{eq:7.4}，}\quad
  \int F_0(\theta)\,\mathrm{d}\theta\leq 2\varepsilon c.
\]
由于 $\varepsilon$ 任意且 $F_0(\theta)$ 非负，我们得出 $F_0(\theta)=0$ 对几乎所有的 $\theta$ 成立，换言之，几乎所有方向上的所有平面与 $E$ 交成平面测度为零的集合.

这类论证也可以反过来用.我们知道（\cref{thm:7.3}）存在在每个方向含有一条直线的零测度平面集，因此若形如 \eqref{eq:7.4} 的估计对 $\mathbb{R}^2$ 上的函数成立（此时 $\theta$ 表示平面内的单位向量），与上述类似的论证将导致矛盾.于是 Besicovitch 集为``\eqref{eq:7.4} 对平面函数成立''这一看似合理的猜想提供了反例.

如\cref{thm:7.13} 及其前的注记所述，$(n,k)$-Besicovitch 集在 $k>1$ 时不存在.有趣的是，上述方法经调整可以在 $k>\frac{1}{2}n$ 时证明这一点，但此外则不然（见 \textcite{falconer1980a}）.这或许暗示：当 $2\leq k\leq\frac{1}{2}n$ 时，$(n,k)$-Besicovitch 集的不存在性是一种内在的几何性质，而非更一般的泛函分析结果的推论.

我们转向球面平均与球的填装问题.若 $f$ 是 $\mathbb{R}^n$ 上的可测函数，令 $F_r(x)$ 表示 $f$ 在以 $x$ 为心、$r$ 为半径的球面上的平均，即 $f$ 在该球面上的积分除以球的 $(n-1)$ 维曲面面积.令
\begin{equation}\label{eq:7.5}
  F(x)=\sup_{r>0}F_r(x).
\end{equation}
\textcite{stein1976} 证明了
\begin{equation}\label{eq:7.6}
  \|F\|_{p}\leq c\|f\|_{p},
\end{equation}
只要 $n/(n-1)<p$ 且 $n\geq 3$.（如果我们改为把 $F_r(x)$ 定义为 $f$ 在以 $x$ 为心、$r$ 为半径的球体上的平均，则 $F(x)$ 就是 Hardy--Littlewood \newterm{maximal-operator}{极大算子}{maximal operator}，而 \eqref{eq:7.6} 就是经典的极大不等式，它对一切 $n$ 与 $1<p<\infty$ 成立，见 \textcite{deguzman1975}.）

我们使用与上面完全相同的论证，但从 \eqref{eq:7.6} 而非 \eqref{eq:7.4} 出发，推出：若 $n\geq 3$ 且 $E$ 是 $\mathbb{R}^n$ 中 $n$ 维勒贝格测度为零的任意子集，则对几乎所有的 $x\in\mathbb{R}^n$，以 $x$ 为心的一切球与 $E$ 交成 $(n-1)$ 维曲面测度为零的集合.

若 \eqref{eq:7.6} 对 $n=2$ 与某个 $p$ 值得证，则将蕴含零测度平面集不可能包含以每点为心的圆，从而解决前文提到的问题.另一方面，若构造出这样的零测度圆填装，则将蕴含 \eqref{eq:7.6} 对一切 $p$ 都不成立，而不只是对 $p\leq 2$ 不成立——后者是目前在 $\mathbb{R}^2$ 中仅知的.

\eqref{eq:7.6} 的证明关键依赖于球面曲率的不为零.对其他``曲率良好''曲面上的平均也有类似结果（见 \textcite{steinwainger1978}），并随之得到填装结果；但对立方体的表面这类结果则不成立，例如.

勒贝格密度定理的许多推广涉及球以外的集合上的平均.一个看似合理的猜想（对连续函数当然成立）或许是：若 $f$ 是平面上的``合理''函数，则
\begin{equation}\label{eq:7.7}
  \inf\Bigl\{\frac{1}{\mathscr{L}^{2}(R)}\int_{R}f\,\mathrm{d}\mathscr{L}^{2}:\ R\ \text{是中心在}\ x\ \text{处且}\ 0<|R|\leq r\ \text{的矩形}\Bigr\}
\end{equation}
当 $r\to 0$ 时对几乎所有的 $x$ 收敛于 $f(x)$，对上确界也有类似结果.然而，Nikodym 集的存在性（在\cref{thm:7.3} 之后讨论）的一个推论是：这类结果对集合的特征函数甚至都不成立.（我们说全体矩形之类不构成一个微分基.）事实上，设 $E$ 是 Nikodym 集，即单位正方形中测度为 $1$、且每一点都位于某条与 $E$ 仅交一点的直线上的子集.由正则性，可找到 $E$ 的闭子集 $F$ 具有同样的可达性质且 $\mathscr{L}^2(F)>0$.令 $f$ 为 $F$ 的特征函数.由于 $F$ 的补是开集，容易看出：取长轴沿``例外直线''的细矩形，\eqref{eq:7.7} 的下确界当 $r\to 0$ 时对一切 $x\in F$ 趋于零.

Besicovitch 集在 \newterm{fourier-multiplier}{Fourier 乘子理论}{Fourier multiplier}中提供了一个重要而惊人的反例.设 $T$ 是 $\mathbb{R}^2$ 上 $p$ 次方可积函数空间上的线性变换，用 Fourier 变换记号定义为
\[
  \widehat{(Tf)}(\xi)=\chi_{\mathrm{B}}(\xi)\hat{f}(\xi),
\]
其中 $\chi_{\mathrm{B}}$ 是单位圆盘的特征函数.已知
\begin{equation}\label{eq:7.8}
  \|Tf\|_{p}\leq c\|f\|_{p}
\end{equation}
对 $p=2$ 成立，且人们曾普遍预期这一不等式对一切 $\frac{4}{3}<p<4$ 也成立.然而 \textcite{fefferman1971} 精妙地利用 Besicovitch 集证明：\eqref{eq:7.8} 当 $p\neq 2$ 时不成立.有兴趣的读者可参阅 Fefferman 行文清晰的文章了解进一步细节.

\textcite{mitjaginnikisin1973} 遵循类似的过程得到：对 $p<2$，存在单位正方形上的 $p$ 次方可积函数 $f$，其\newterm{fourier-series}{Fourier 级数的部分和}{Fourier series}几乎处处无界，即对几乎所有的 $(x,y)\in\mathbb{R}^2$，
\[
  \varlimsup_{r\to\infty}\Bigl|\sum_{k^{2}+j^{2}<r}a_{kj}\exp\bigl(2\pi\mathrm{i}(xk+yj)\bigr)\Bigr|=\infty,
\]
其中 $\{a_{kj}\}$ 是 $f$ 的二维 Fourier 系数.这与一维情形形成鲜明对比：Carleson--Hunt 定理断言，若 $f$ 对某个 $1<p<\infty$ 是 $p$ 次方可积的，则 $f$ 的 Fourier 级数几乎处处逐点收敛.

最后我们提到在\newterm{functional-analysis}{泛函分析}{functional analysis}中的一个应用：$T$ 是 Hilbert 空间上的 \newterm{hypernormal-operator}{超正规算子}{hypernormal operator}（即 $T^{*}T-TT^{*}\geq 0$ 且 $T$ 没有非平凡的约化空间使其上 $T$ 正规），当且仅当 $T$ 的谱在其每一点的任何邻域中都具有正测度.\textcite{putnam1974} 因此用 Nikodym 集证明了一个正面结果，而不仅仅是提供反例.

要更深入地探索调和分析的相关领域，读者可参阅 \textcite{steinwainger1978} 的综述文章、由 \textcite{waingerweiss1979} 编辑的美国数学会调和分析会议文集，以及 de Guzmán（\citeyear{deguzman1975,deguzman1981}） 的书.几何测度论与调和分析看来必将继续相互产生巨大的影响.

\begin{exercises}
\item 证明存在 $\mathbb{R}^n$（$n\geq 2$）的子集 $E$，在其中有 $\mathscr{L}^n(E)$ 任意小的测度，且在其内部一条单位长线段可以被挪动到任意方向.
\item 设 $S$ 是 $[0,\pi)$ 的 Borel 子集.设 $F$ 是 $\mathbb{R}^2$ 中在 $S$ 的每个方向都含有一条直线的子集.用对偶性证明 $\dim F\geq 1+\dim S$.
\item 证明存在 $\mathbb{R}^2$ 的 $\mathscr{L}^2$-测度为零的子集，包含过 $x$ 轴每一点的一条互不相同的直线.
\item 证明若 $E$ 是 $\mathbb{R}^2$ 中 $\dim E>1$ 的 Borel 子集，则 $E$ 中各点的极线之集（关于任何固定的圆）覆盖一个平面勒贝格测度为无穷的集合.
\item 若 $f$ 是平面可测函数，令 $F_r(x)$ 表示 $f$ 在以 $x$ 为心、$r$ 为半径的圆周上的平均.可以证明（\citeyear{wainger1979}）若 $E$ 是 Cantor 集，则对某个常数 $c$ 有
\[
  \bigl\|\sup_{r\in E}F_{1+r}(x)\bigr\|_{2}\leq c\|f(x)\|_{2}.
\]
进而证明：勒贝格测度为零的平面集不可能在 $\mathbb{R}^2$ 的每一点处包含一个半径为 $1+r$（对某个 $r\in E$）的圆.
\end{exercises}
